% SpecBridge PRD - AI for Product Managers capstone (IIIT Hyderabad)
% Compile with pdflatex (TeX Live, MiKTeX or Overleaf): pdflatex SpecBridge_AI_PRD.tex
\documentclass[10pt,a4paper]{article}

\usepackage[a4paper,left=14mm,right=14mm,top=12mm,bottom=13mm]{geometry}
\usepackage[T1]{fontenc}
\usepackage[utf8]{inputenc}
\usepackage{lmodern}
\usepackage[scaled=0.94]{helvet}
\renewcommand{\familydefault}{\sfdefault}
\usepackage{microtype}
\usepackage{xcolor}
\usepackage{array,tabularx,booktabs}
\usepackage{enumitem}
\usepackage[most]{tcolorbox}
\usepackage[hidelinks]{hyperref}

% ---- colours (same palette as the web version) ------------------------------
\definecolor{ink}{HTML}{1B2230}
\definecolor{muted}{HTML}{5B6575}
\definecolor{accent}{HTML}{2B57C7}
\definecolor{accentsoft}{HTML}{E8EEFB}
\definecolor{rulegrey}{HTML}{DFE3EA}
\definecolor{good}{HTML}{1D7A57}
\definecolor{warn}{HTML}{A8571C}
\definecolor{bad}{HTML}{B03A3A}

% ---- page and text ---------------------------------------------------------
\pagestyle{empty}
\setlength{\parindent}{0pt}
\setlength{\parskip}{3.5pt}
\setlist{nosep,leftmargin=1.2em,itemsep=2pt}
\renewcommand{\arraystretch}{1.25}
\setlength{\tabcolsep}{4pt}
\arrayrulecolor{rulegrey}
\newcolumntype{L}{>{\raggedright\arraybackslash}X}
\newcolumntype{P}[1]{>{\raggedright\arraybackslash}p{#1}}

% ---- small building blocks ---------------------------------------------------
\newcommand{\prdsec}[2]{\par\vspace{7pt}{\large\bfseries{\color{accent}#1}\hspace{0.55em}#2}\par\vspace{3pt}}
\newcommand{\eyebrow}[1]{{\footnotesize\bfseries\color{accent}\textls[120]{\MakeUppercase{#1}}}}
\newcommand{\boxlabel}[1]{{\scriptsize\bfseries\color{accent}\textls[100]{\MakeUppercase{#1}}}\par\vspace{2pt}}
\newcommand{\thead}[1]{{\scriptsize\bfseries\color{muted}\textls[80]{\MakeUppercase{#1}}}}
\newcommand{\note}[1]{{\small\color{muted}#1}}
\newcommand{\assume}[1]{{\color{warn}#1}}
\newcommand{\sev}[2]{\tcbox[on line,colframe=#1,colback=white,boxrule=0.5pt,arc=1pt,
  left=2pt,right=2pt,top=0.5pt,bottom=0.5pt]{\scriptsize\bfseries\ttfamily\color{#1}#2}}

\tcbset{
  jtbd/.style={colback=accentsoft,colframe=accentsoft,boxrule=0pt,arc=2pt,
    left=7pt,right=7pt,top=6pt,bottom=6pt},
  card/.style={colback=white,colframe=rulegrey,boxrule=0.5pt,arc=2pt,
    left=7pt,right=7pt,top=6pt,bottom=6pt},
  ticket/.style={card,enhanced,borderline west={2pt}{0pt}{accent}},
}

\begin{document}
\color{ink}

% ============================================================================
% Header
% ============================================================================
\eyebrow{AI for Product Managers \textperiodcentered{} IIIT Hyderabad \textperiodcentered{} Capstone project \textperiodcentered{} PRD}

\vspace{2pt}
{\fontsize{19}{23}\selectfont\bfseries SpecBridge: Jira tickets to codebase-grounded UAT test cases\par}

\vspace{3pt}
{\small\textbf{Team} {\color{muted}Unmanned Rooks: Hussain Lohawala, Qasim Naik, Rishabh Singh, Dhawal}
\hspace{1.5em}\textbf{Version} {\color{muted}v0.1 \textperiodcentered{} 28 Oct 2026}\par}

\vspace{4pt}
Today a PM sees the user journey only after engineers have changed 3--5 repositories, and if it is not
what the PM meant, the work goes round again. SpecBridge reads the team's real code and shows the PM
that journey first, as UAT test cases (steps, expected results, affected screens, open questions), so
PM and engineers agree on it before anything is built.

\vspace{3pt}
{\color{ink}\rule{\linewidth}{1.2pt}}

% ============================================================================
\prdsec{1}{User and job-to-be-done}
\begin{tcbraster}[raster columns=2,raster equal height,raster column skip=10pt]
\begin{tcolorbox}[jtbd]
\boxlabel{User}
A product manager who hands UI tasks to frontend and backend engineers on a product whose
frontend is split across several repositories.
\end{tcolorbox}
\begin{tcolorbox}[jtbd]
\boxlabel{Job-to-be-done}
When I hand over a UI task, I want to see and agree the exact user journey, including errors and
edge cases, before engineers build it, so the first demo is what I meant.
\end{tcolorbox}
\end{tcbraster}
\note{Secondary users: the frontend engineer, who builds from an agreed journey instead of guessing,
and QA, who later runs the same cases.}

% ============================================================================
\noindent
\begin{minipage}[t]{0.485\linewidth}
\prdsec{2}{Problem}
\textbf{How it works today:}
\begin{enumerate}
  \item The PM gives a UI task to the frontend and backend engineers.
  \item The frontend engineer translates it into changes across 3--5 repositories.
  \item Only when the work is done does the engineer demo the user journey to the PM.
  \item If it is not what the PM meant, the cycle repeats, sometimes several times.
\end{enumerate}
\vspace{3pt}
\textbf{Why:} the journey is never made concrete and agreed before building. The PM writes in product
language (``let users finalize a bill'') and does not know which screens and flows across the repos
the change touches, so mismatches, error states and edge cases surface at the demo, the most
expensive moment to find them.

\textbf{Why now:} current LLMs can read code and plain language together, and a ticket costs cents
to process.
\end{minipage}\hfill
\begin{minipage}[t]{0.485\linewidth}
\prdsec{3}{Approach}
\begin{tcbitemize}[raster columns=3,raster equal height,raster column skip=4pt,raster row skip=4pt,
  colback=white,colframe=rulegrey,boxrule=0.5pt,arc=1pt,left=3pt,right=3pt,top=2pt,bottom=2pt,
  fontupper=\footnotesize\ttfamily]
\tcbitem {\color{muted}1}\ Ticket URL
\tcbitem[colframe=accent,coltext=accent] {\color{muted}2}\ Understand
\tcbitem {\color{muted}3}\ Orchestrator
\tcbitem[colframe=accent,coltext=accent] {\color{muted}4}\ Criterion agents $\times$N
\tcbitem[colframe=accent,coltext=accent] {\color{muted}5}\ Merge + review
\tcbitem {\color{muted}6}\ Confluence page
\end{tcbitemize}
\vspace{2pt}
\begin{itemize}
  \item \textbf{Understand} the ticket: criteria (stated and inferred), screens, questions, risks.
        Unclear or unsafe tickets get questions, not invented cases.
  \item \textbf{Agents per acceptance criterion} search a code knowledge base (keywords + local
        embeddings + links between files across repos) with read-only tools, then write cases that
        cite the code they read.
  \item \textbf{An independent reviewer} checks coverage and that every cited screen exists; weak
        criteria go back for another round.
  \item \textbf{Model:} Gemini 3.8 Flash via OpenRouter; the model can be swapped through
        configuration. Built with LangGraph.
\end{itemize}
\end{minipage}

% ============================================================================
\prdsec{4}{Example}
\begin{tcbraster}[raster columns=2,raster column skip=10pt,raster valign=top,
  raster force size=false]
\begin{tcolorbox}[ticket,width=0.36\linewidth]
\boxlabel{The ticket, as the PM writes it}
{\bfseries ``Start visit'' does nothing from patient search\par}
\vspace{3pt}
When the front-desk user searches for a patient and clicks \textbf{Start visit} on a result, nothing
happens. Expected: the visit form opens for that patient and the visit can be started.
\end{tcolorbox}
\begin{tcolorbox}[card,width=0.62\linewidth]
\boxlabel{What SpecBridge returns {\normalfont\color{muted}(illustrative)}}
\note{\textbf{Affected screens:} patient search results, where the button sits, and the start-visit
form. The button is built in the patient-chart repo and shown in the patient-management search
screen, so two repos are involved.}

\small
\begin{tabularx}{\linewidth}{@{}P{12mm}LL@{}}
\toprule
\thead{Case} & \thead{Steps} & \thead{Expected result}\\
\midrule
\texttt{TC-01}\newline{\color{muted}\scriptsize happy} & Log in as front desk. Search a patient with no
  active visit. Click \textbf{Start visit} on the result. & The start-visit form opens for that
  patient; after saving, the visit shows as active.\\
\midrule
\texttt{TC-02}\newline{\color{muted}\scriptsize error} & Same steps, while the server rejects the
  request. & An ``Error starting visit'' message appears; no visit is created.\\
\bottomrule
\end{tabularx}

\note{\textbf{Question for the PM:} if the patient already has an active visit, should the button be
hidden or open that visit?}
\end{tcolorbox}
\end{tcbraster}
\note{The PM reviews this before any code is written. A mismatch, like a different answer to the
question above, is caught now instead of at the demo.}

% ============================================================================
\newpage
\prdsec{5}{Build vs rules}
\noindent
\begin{minipage}[t]{0.485\linewidth}
\textbf{Why rules are not enough}
\begin{itemize}
  \item \textbf{Free-text input.} Tickets come in any style, from formal criteria to informal
        Hinglish. Templates enforce structure, not correctness.
  \item \textbf{Language-to-code mapping is semantic.} ``Finalize a bill'' must be matched to the
        right code among 2,322 files in 5 repos. On our golden tickets with no file names, keyword
        search plus code links finds a changed file in the top 10 for 74\% of tickets; adding AI
        search by meaning raises it to 79\%. The agent then searches further on its own.
  \item \textbf{Reasoning, not lookup.} Inferring unstated criteria, negative paths and
        contradictions in a ticket needs judgment.
\end{itemize}
\vspace{4pt}
\textbf{Where we keep rules (deterministic guardrails)}\par
Team-only ticket intake, secret redaction, read-only AI tools, verification of every cited file and
screen, publishing done only by code, cost and time caps.

\vspace{4pt}
\textbf{Why AI is worth it}\par
An AI draft costs about \$0.02--0.07 and 1--2 minutes per ticket (estimate from our measured token
sizes), against \assume{an assumed 30--60 minutes} of PM time per ticket today, and it brings code
knowledge the PM lacks. \assume{The time saving is an assumption we will measure in customer
sessions.}
\end{minipage}\hfill
\begin{minipage}[t]{0.485\linewidth}
\textbf{Accuracy needed and what errors cost}

\small
\begin{tabularx}{\linewidth}{@{}LP{27mm}P{17mm}@{}}
\toprule
\thead{Error} & \thead{Cost} & \thead{Tolerance}\\
\midrule
Leaks a secret, or obeys instructions hidden in a ticket & Security incident & \sev{bad}{ZERO}\\
\midrule
Invents a screen or flow that does not exist & Wasted PM/dev time; trust lost & \sev{warn}{LOW}\\
\midrule
Misses a case (e.g.\ an error state) & Bug found later; rework (PM review catches most) & \sev{good}{SOME}\\
\midrule
Awkward wording or a duplicate case & Minor edit & \sev{good}{SOME}\\
\bottomrule
\end{tabularx}

\note{Output is a draft a human approves, so it need not be perfect; grounding and safety must be
near-perfect.}
\end{minipage}

% ============================================================================
\prdsec{6}{Risks and mitigations}
\small
\begin{tabularx}{\linewidth}{@{}P{0.31\linewidth}L@{}}
\toprule
\thead{Risk} & \thead{Mitigation}\\
\midrule
Hallucinated screens, steps or code references & Every case must cite code it read; the reviewer
  verifies each cited file, line and screen name against the code.\\
\midrule
Prompt injection or secrets inside tickets & Ticket text treated as data; secrets removed before the
  AI sees them; AI has read-only tools; only code publishes.\\
\midrule
Vague or contradictory tickets & A gate turns them into questions for the PM instead of guessed test
  cases.\\
\midrule
Search misses the right code (79\% hit rate on fair tickets today) & Agents search again per
  criterion; misses are logged and fixed in each improvement round.\\
\midrule
The build still drifts from the agreed journey & The approved page is the engineer's direction and
  the checklist for the demo, so drift is caught step by step.\\
\midrule
Overwriting the PM's edits on re-runs & Updates never overwrite human edits; new drafts go in as a
  page comment.\\
\midrule
Cost or latency runaway; free-tier model capacity & Caps of \$0.25 and 2 review rounds per ticket;
  fallback from free to paid capacity.\\
\midrule
Our evaluation flatters us & Real tickets fixed after a frozen code version, answer key from the
  real fixes, links to the fix removed; 18 tickets that name code files scored separately;
  6 adversarial tickets.\\
\bottomrule
\end{tabularx}
\normalsize

% ============================================================================
\noindent
\begin{minipage}[t]{0.485\linewidth}
\prdsec{7}{Success metric}
\textbf{Outcome we want:} fewer demo-and-rework rounds per UI task.
\note{Baseline: rounds per task today, asked in customer interviews.}

\textbf{Primary (measurable in the capstone):} share of generated test cases the PM keeps without
edits, i.e.\ the journey was right before building.
\note{Measured in 5+ sessions with real PMs, engineers or QA (not classmates), with verbatims.}

\textbf{Also tracked:} mismatches the page surfaced before build (per PM), engineer clarity rating
(1--5), tokens, cost and time per ticket.

\note{\textbf{Test bed:} OpenMRS 3, 5 real connected React repos frozen at 31 Mar 2026; golden set
of 63 real tickets fixed after that date (answer key from the real code changes) + 6 hand-written
adversarial tickets. The team picks and labels the final 30--50.}
\end{minipage}\hfill
\begin{minipage}[t]{0.485\linewidth}
\prdsec{8}{Ship gate {\normalsize\normalfont\color{muted}(proposed, team to confirm)}}
\small
\begin{tabularx}{\linewidth}{@{}L>{\raggedleft\arraybackslash\ttfamily}p{22mm}@{}}
\toprule
\thead{Measure (golden set)} & \thead{Gate}\\
\midrule
Acceptance criteria covered by at least one case & $\geq$ 90\%\\
\midrule
Cited files and screens verified in code (grounding) & $\geq$ 95\%\\
\midrule
Adversarial tickets handled safely & 6 / 6\\
\midrule
Human rubric, 2 raters (accuracy, grounding, usefulness) & $\geq$ 4 / 5\\
\midrule
LLM judge agrees with humans on spot-checks & $\geq$ 80\%\\
\midrule
Cost / time per ticket & $\leq$ \$0.10 / 3 min\\
\bottomrule
\end{tabularx}
\normalsize
\end{minipage}

% ============================================================================
\vfill
{\color{rulegrey}\rule{\linewidth}{0.5pt}}
\note{\footnotesize Out of scope: writing or changing code, automated tests, backend design, running
the app, auto-approving pages. Sources: team proposal (Oct 2026), project build log and measured
search scores (Step 4, 27 Oct 2026).}

\end{document}
